% NOTES ONLY, to be written. Trimmed 2026-11-07 from thesis/methodology/07_curvature_fpca.tex
% (full draft in git, commit e4eb857). Decision 2026-11-04: main model uses mean curvature; FPCA
% shape scores only in pLAD and mLAD, each its own descriptive cloglog model, not pooled, not part
% of the LR test. Theory: sec:tech-fpca, sec:tech-fpca-grid, sec:tech-fpca-registration.
% Settings: CGPS_data/turning_profile_shapes.py and segment_k_budget.py, neither in thesis_final;
% turning_profile_shapes still runs on the Medis trees.
%
% The section should contain:
% - Why: mean curvature gives the same value to one sharp bend and to even turning; FPCA of the
%   profile describes where along the segment the turning lies. Secondary, descriptive, outside the
%   test of the geometry block (sec:method-models). Restricted to pLAD and mLAD for events ([..]).
% - Profile (eq:method-profile) on 50 equally spaced points of normalised position, linear
%   interpolation; FPCA per segment by SVD of the centred participant x grid matrix; no registration,
%   so a component can also be a shift of a bend.
% - K fixed before outcome linkage: smaller of [rule] and the Riley budget of the segment less the
%   base. Open: which rule (80/90% variance, Horn, circular-shift null, bootstrap |cos| >= 0.9,
%   5-fold CV error).
% - Per segment a cloglog model with the base and K scores, one row per participant, estimates with
%   intervals, not tested. Open: the base (1 df, recalibrated SCORE2, in segment_k_budget.py) and
%   the interval method.
\section{Secondary Analysis of Curvature Shape}
\label{sec:method-fpca}

\todonote{To be written: FPCA of the curvature profile in the proximal and mid LAD as a descriptive secondary analysis.}
